\section{Conclusion}
\label{sec:conclusion}

Table~\ref{tab:synthesis}, at the head of \S\ref{sec:findings}, states what this corpus reports, what it
stops short of, what limits each conclusion and the experiment that would test it. How firmly each of its
five conclusions is supported differs, and is stated here rather than compressed into a grade.
\emph{Availability against use} rests on the widest spread of distinct first authors in the corpus
--- the decodability rows draw on \Fdecgrpmin{} to \Fdecgrpmax{} distinct first authors over four to seven
modalities --- but on records that almost never test selectivity, so the firm part is that a readout
recovers the target, not that the model uses it. \emph{Acquisition and site structure} is the property
most repeatedly reported in the corpus, by many groups and in several modalities, but not the
best controlled: a comparison model appears in a minority of the records carrying F3a and F3c, and one
modality dominates. \emph{Meaning of discovered features} rests on \Nsae{} dictionaries of which
\Nsaeval{} report a semantic validation and \Nsaestab{} any reproducibility check, so the negative half of
the conclusion, that cross-seed identity is unestablished, is firmer than the positive half.
\emph{Medical against general pre-training} is the weakest: the comparisons are uncontrolled by
construction, so the conclusion is an absence of identification rather than a finding. \emph{Where to
measure} rests on three heterogeneous studies and one successive-locus probe and is offered as a
hypothesis. Two qualifications apply to all five: \Pctpreprint\% of the corpus is preprints,
and no formal risk-of-bias instrument was applied, so the appraisal is the reporting audit of
Fig.~\ref{fig:findings} rather than a certainty grade. Three consequences for a reader who studies or
deploys a medical foundation model follow.

\emph{Breadth is not control.} The checks these studies run are not interchangeable, so they are counted
in \Vnclasses{} classes, each stated with what it licenses and what it does not (Table~\ref{tab:vcov}).
A statement that lacks a class is narrower, not refuted. A probe establishes that a readout recovers
the target, and no observational control alone makes it establish use; a control task adds selectivity and a random encoder
adds attribution to training, each a different, still weaker statement than use. A dictionary feature is
a clinical concept only after clinician or structured-label validation, however it was named, and a
steering effect is direction-selective only against matched controls; no medical study in this corpus
meets the criterion of \S\ref{sec:methods} for a concept-specific causal contribution.

\emph{What to measure.} Recovery is reported broadly and use is barely measured, and reported broadly is
not the same as established: \Fsensrefemptyn{} rows lose every record when the analysis is restricted
to peer-reviewed studies, only \Nminvalid{} of the \Ncore{} studies state an independent split, evaluate
on unused data and report uncertainty, and capacity and representation-attribution controls on probes
are rare, with no validated image-domain analogue of the linguistic randomized control task yet
established. A new audit should therefore not add another probe: it should pair one with a
matched intervention or an erasure at the same locus and report both.

\emph{Where to measure.} The readout, whether pooling, depth or the direction chosen, can change what
appears recoverable, and the connector output and text encoder of multimodal models remain almost
unmeasured as loci of their own; a locus should be named before a model is praised. Whether medical
pre-training itself changes any of this is unidentified in the same-data comparisons that exist.
